% 자동 생성 — scripts/build_tables.py. 손으로 고치지 않는다.
% 출처: output/results/21_b_group.json : 사다리
\begin{table}[htbp]
  \centering
  \caption{2021.4 B그룹 --- 기준점을 바꿔도 $-41.7 \sim -44.9\%$ 다.}
  \label{tab:ladder-b}
  \begin{tabular}{lrrrr}
    \toprule
    \textbf{기준점} & \textbf{동} & \textbf{효과 (\%)} & \textbf{WCR $p$} & \textbf{무작위화 $p$} \\
    \midrule
    ① 서울 전역 비지정 & 182 & -44.7 & 0.057 & $<0.001$ \\
    ② ① + 가격대 공통지지 & 39 & -41.7 & 0.063 & 0.002 \\
    ③ 오염 미제외 & 233 & -43.8 & 0.071 & $<0.001$ \\
    ④ A그룹 맞닿은 동 포함 & 188 & -44.9 & 0.063 & $<0.001$ \\
    \bottomrule
  \end{tabular}

  \begin{minipage}{\linewidth}\footnotesize
  \textit{주.} 목동이 속한 양천구에 쓸 수 있는 비지정 동이 0개라 `같은 구 안' 통제군을 주 사양으로 쓸 수 없었다.
  \end{minipage}
\end{table}
